% ECONOMICS_PROBLEM: {"id": "H22-617", "title": "Tax salience and behavioral incidence", "jel_code": "H22", "source_id": "OP-0046"}
% !TeX program = xelatex
\documentclass[11pt,a4paper]{article}
\usepackage[margin=23mm]{geometry}
\usepackage{fontspec}
\setmainfont{Latin Modern Roman}
\usepackage{amsmath,amssymb,mathtools}
\usepackage{xcolor,enumitem,booktabs,longtable,array,xurl,hyperref}
\definecolor{muted}{HTML}{53636C}
\hypersetup{colorlinks=true,urlcolor=blue,linkcolor=blue}
\setlength{\parindent}{0pt}
\setlength{\parskip}{5pt}
\newcommand{\R}{\mathbb R}
\newcommand{\E}{\mathbb E}
\newcommand{\N}{\mathbb N}
\newcommand{\ind}{\mathbf 1}
\newcommand{\Var}{\operatorname{Var}}
\newcommand{\Cov}{\operatorname{Cov}}
\newcommand{\supp}{\operatorname{supp}}
\DeclareMathOperator*{\argmin}{arg\,min}
\DeclareMathOperator*{\argmax}{arg\,max}
\newcommand{\entrymeta}[3]{{\sffamily\small\color{muted}\textbf{\texttt{#1}}\quad|\quad #2\par #3\par}}
\newcommand{\Problemref}[1]{\texttt{#1}}
\newcommand{\JELref}[2]{\texttt{#2} (JEL \texttt{#1})}
\begin{document}
\section*{Tax salience and behavioral incidence}
\textbf{Problem ID:} \texttt{H22-617}\quad \textbf{JEL:} \texttt{H22}\par
% BEGIN ECONOMICS STATEMENT
\label{problem:OP-0046}
\label{atlas:prob:P6}
\noindent\textbf{Original atlas identifier:} P6; entry 46. \textbf{Field:} Public finance and taxation. \textbf{Original tractability label:} Medium.

\noindent\textbf{Classification:} Parameterized theoretical/empirical research agenda; not a certified unsolved theorem.

\subsection*{Definitions and assumptions}
Fix a tax liability schedule and a market with true tax-inclusive price $q$, posted pretax price $p$, and consumer attention state $a_{it}\in\mathcal A$. A presentation intervention $d\in\mathcal D$ changes display or payment timing but holds the legal liability and feasible consumption set fixed. A behavioral model $m$ maps information and experience into $a_{it}$ and demand $x_{it}$, including learning, memory, and any effort cost of attention. Actual welfare is evaluated using a prespecified experienced-utility or informed-choice criterion $U_i$, not automatically the utility rationalizing inattentive demand. Let $EV_i(d;m)$ be its money-metric equivalent variation, $R(d;m)$ government revenue, and $C(d;m)$ real presentation/administrative cost. Fix weights $w_i\geq0$ and public-fund value $\lambda\geq0$.

\subsection*{Formal research question}
\emph{Editorial formulation: the mathematical target below is not a theorem quoted from the references.}

For periods $t=0,\ldots,T$, identify dynamic treatment effects
\[
 \Delta_t=\mathbb E[x_{it}(d_1)-x_{it}(d_0)],
\]
where the two potential demands correspond to economically equivalent liability schedules and prespecified intervention histories. Recover or bound $\mathbb E[w_iEV_i(d;m)]+\lambda[R(d;m)-C(d;m)]$ under each declared normative criterion. Determine whether attention responses converge, persist, or vary with repeated exposure, and which experiments distinguish learning from habit formation or compositional attrition. The causal equality of tax liabilities is not an assumption that attention states or equilibrium producer prices are equal; specify the price-setting model when presentations change demand. Characterize policy rankings over models observationally equivalent in purchase data.

\subsection*{Interpretation and scope}
A salience effect is a difference between interventions, not proof of welfare loss by itself. In particular, demand responses cannot identify welfare without a benchmark for experienced utility, beliefs, and the cost of becoming informed. Tax-inclusive displays can reveal information and simultaneously alter search or framing, so a factorial design may be needed to distinguish mechanisms. Long follow-up requires outcomes for the initial assigned sample; analyzing continuing purchasers alone can manufacture apparent learning. Statutory and economic incidence need not coincide under limited attention, but legal remittance changes must be separated from the display intervention. The substantive question concerns persistence and mechanism-dependent welfare, with the short-run existence of salience effects already addressed by the references.

\subsection*{Source audit and status}
[S1]--[S3] are identifiable publications with existing theory and evidence. The original question is therefore not an unanswered yes/no question about whether visibility ever matters. The editorial extension concerns dynamic attention and welfare in a specified context, without a literature-wide 2026 openness certification.

\subsection*{References}
\begin{enumerate}[label={[S\arabic*]},leftmargin=*,itemsep=0.6em]
\item Raj Chetty, Adam Looney, and Kory Kroft (2009). Salience and Taxation: Theory and Evidence. American Economic Review 99(4), 1145--1177. DOI: 10.1257/aer.99.4.1145.
\url{https://pubs.aeaweb.org/doi/abs/10.1257/aer.99.4.1145}
\newline\emph{Locator and support limits:} Abstract and experiments: salience changes behavior and tax incidence; not identification of long-run learning.
\item Amy Finkelstein (2009). E-ZTax: Tax Salience and Tax Rates. The Quarterly Journal of Economics 124(3), 969--1010. DOI: 10.1162/qjec.2009.124.3.969.
\url{https://academic.oup.com/qje/article-abstract/124/3/969/1905173}
\newline\emph{Locator and support limits:} Abstract and toll analysis: salience and policy rates in one institutional setting.
\item Dmitry Taubinsky and Alex Rees-Jones (2018). Attention Variation and Welfare: Theory and Evidence from a Tax Salience Experiment. The Review of Economic Studies 85(4), 2462--2496. DOI: 10.1093/restud/rdx069.
\url{https://academic.oup.com/restud/article-abstract/85/4/2462/4665811}
\newline\emph{Locator and support limits:} Abstract and model: heterogeneous attention and welfare; normative interpretation requires the stated preference benchmark.
\end{enumerate}

\subsection*{Common conventions: Source mathematical conventions}

Unless an entry states otherwise, agent and firm sets are finite. State and action spaces are measurable, strategies and outcome maps are measurable, probability laws are specified, and expectations exist for the targets under discussion. Infinite-horizon discount factors lie in $(0,1)$ and discounted objectives are finite. Norms, tolerances and complexity input sizes must be specified for computational comparisons. Constants or primitives that appear locally are inputs, not empirical facts asserted by this catalogue.

\textbf{Identification.} A statistical model $m\in\mathcal M$ implies an observable law $P_m$ and target $\theta(m)$. At law $P$, its identified set is
\[
\Theta(P;\mathcal M)=\{\theta(m):m\in\mathcal M,\ P_m=P\}.
\]
Point identification holds when this set is a singleton, or equivalently when $P_{m_1}=P_{m_2}$ implies $\theta(m_1)=\theta(m_2)$ for all compatible models. Sharp bounds mean bounds attained, or approached when the boundary is not attained, by compatible models. Writing this set is a definition, not a solution: the substantive task is to establish informative restrictions and derive or estimate the set. An unrestricted-confounding model may give only trivial bounds.

\textbf{Inference.} Identification, estimability and valid uncertainty quantification are separate questions. Fix $\alpha\in(0,1)$. A confidence procedure $C_n$ over a declared law class $\mathcal P$ has asymptotic uniform coverage at level $1-\alpha$ when
\[
\liminf_{n\to\infty}\inf_{P\in\mathcal P}\Pr_P\{\theta(P)\in C_n\}\geq1-\alpha.
\]
For set-identified targets the coverage event must be defined separately, for example coverage of the full identified set. If an entry uses identification-indexed models instead of uniquely indexed laws, the same coverage convention is applied over those models. Pointwise limits do not establish uniform validity, and asymptotic validity does not guarantee finite-sample precision. Sample sizes, dependence structures and weak-identification sequences are part of each application.

\textbf{Counterfactuals.} $Y_i(a)$ is a potential outcome under a specified intervention, including its timing, implementation and versions. It is not $Y_i$ conditional on voluntarily choosing $a$. A full intervention vector permits interference; shorthand scalar treatment notation does not silently impose no interference. Assignment, selection, attrition, anticipation and measurement must be specified. A claim about an instrument's local effect is not automatically a population policy effect.

\textbf{Economic equilibrium.} A counterfactual model includes best responses, constraints, feasibility, market clearing, state transitions and expectations consistent with the stated information. When several equilibria exist, either a declared selector is an input or the target is set-valued. Comparisons across policy regimes must use the stated selection convention. Reduced-form relationships are not presumed invariant after an equilibrium-changing intervention.

\textbf{Optimization and welfare.} Optimization is over the stated feasible set. If attainment is not established, $\sup$ or $\inf$ and $\varepsilon$-optimal choices replace an asserted maximizer or minimizer. For example, with value $V=\sup_{a\in\mathcal A}W(a)<\infty$, an $\varepsilon$-optimal set is $\{a\in\mathcal A:W(a)\geq V-\varepsilon\}$ for $\varepsilon>0$. Benefits, costs, risk and health or environmental outcomes can be aggregated only using declared units and conversion weights. Interpersonal utility weights, the treatment of fiscal transfers and foreign profits, discounting, and the behavioral welfare criterion are explicit inputs. A comparative-static derivative additionally requires existence and appropriate differentiability of the selected counterfactual.

\textbf{Sources.} Local labels [S1], [S2], and so forth restart in every question. Locators identify the relevant abstract, model, section or result at the level actually checked; a bibliographic or abstract-level verification is not represented as inspection of an entire proof. Working papers are distinguished from later publications and changing versions. Source summaries are paraphrased. All URL checks and editorial formulations refer to the audit date stated above.
% END ECONOMICS STATEMENT
\end{document}
